\clearpage
\VerdierPhysicalPageAnchor{0146}%
\begingroup
\small
\setlength{\parskip}{0pt}%
\setlength{\emergencystretch}{1.5em}%
\setlength{\abovedisplayskip}{3pt plus 1pt minus 1pt}%
\setlength{\belowdisplayskip}{3pt plus 1pt minus 1pt}%
\sloppy

\begingroup\itshape
\VerdierSemanticAnchor{ch2-2-3-3-f}
\noindent f) Un objet $Z$ de $\mathcal D$ appartient à
$Q^{-1}(\operatorname{Lib}_R(Q))$ si et seulement s'il existe un objet
libre à droite $Y$ et un morphisme $s\colon Z\longrightarrow Y$, où $s$
appartient à $\widehat{S_{\mathcal D}(\mathcal B)}$, le saturé du système
$S_{\mathcal D}(\mathcal B)$ (2.2.4, (c)).

\VerdierSemanticAnchor{ch2-2-3-3-g}
\par\noindent g) Le foncteur d'inclusion
$L_R(Q)\subset Q^{-1}(\operatorname{Lib}_R(Q))$ admet un adjoint à gauche
\[
  \mathcal R\colon Q^{-1}(\operatorname{Lib}_R(Q))\longrightarrow L_R(Q).
\]

\VerdierSemanticAnchor{ch2-2-3-3-adjunction-diagram}
\begin{center}
\begin{tikzcd}[column sep=3.2em,row sep=3.1em]
L_R(Q)
  \arrow[r,shift left=.7ex]
  \arrow[dr,"{Q|_{L_R(Q)}}"']
& Q^{-1}(\operatorname{Lib}_R(Q))
  \arrow[l,dashed,bend right=28,"{\mathcal R}"']
  \arrow[r,hook]
  \arrow[d,"{Q_1}"]
& \mathcal D \arrow[d,"Q"] \\
& \operatorname{Lib}_R(Q)
  \arrow[ul,dashed,bend left=10,"{\Psi}"']
  \arrow[r,hook]
& \mathcal D/\mathcal B
\end{tikzcd}
\end{center}
\endgroup

\VerdierSemanticAnchor{ch2-2-3-3-adjunction-explanation}
Soit $\Psi\colon\operatorname{Lib}_R(Q)\longrightarrow L_R(Q)$ un foncteur
quasi-inverse à $Q|_{L_R(Q)}$ (cf. c) et e)). On peut prendre
$\mathcal R=\Psi Q_1$, où on désigne par
$Q_1\colon Q^{-1}(\operatorname{Lib}_R(Q))\longrightarrow
\operatorname{Lib}_R(Q)$ le foncteur induit par $Q$. Le foncteur $\Psi$ est
adjoint à droite au foncteur $Q_1$.

\VerdierSemanticAnchor{ch2-2-3-4}
\medskip
\noindent\textbf{2.3.4. Démonstration de la proposition 2.3.3 :}\quad
\VerdierSemanticAnchor{ch2-2-3-4-a-proof}
Démontrons a). Tout d'abord, si $Y$ est $Q$-libre à droite, le foncteur
$X\longmapsto\operatorname{Hom}_{\mathcal D}(X,Y)$ est isomorphe au
foncteur
$X\longmapsto\operatorname{Hom}_{\mathcal D/\mathcal B}(Q(X),Z)$, pour un
objet $Z$ convenable de $\mathcal D/\mathcal B$. Par suite, le foncteur
$X\longmapsto\operatorname{Hom}_{\mathcal D}(X,Y)$ transforme les
morphismes de $S_{\mathcal D}(\mathcal B)$ en isomorphismes. Ceci démontre
l'implication i) $\Rightarrow$ ii). L'implication ii) $\Rightarrow$ iii)
résulte de ce que le morphisme $B\longrightarrow0$ est un morphisme de
$S_{\mathcal D}(\mathcal B)$. L'implication iii) $\Rightarrow$ iv) résulte
de la définition de $S_{\mathcal D}(\mathcal B)$ et de (1.2.6).
L'implication iv) $\Rightarrow$ v) résulte de l'isomorphisme (2.2.4) :

\VerdierSemanticAnchor{ch2-2-3-4-colimit-isomorphism}
\[
  \operatorname{Hom}_{\mathcal D/\mathcal B}(Q(X),Q(Y))
  \xrightarrow{\sim}
  \varinjlim_{Y\backslash S_{\mathcal D}(\mathcal B)}
  \operatorname{Hom}_{\mathcal D}(X,\mathord\cdot)\, .
\]

\VerdierSemanticAnchor{ch2-2-3-4-a-proof-complete}
Enfin, l'implication v) $\Rightarrow$ i) est évidente : $Y$ est libre à
droite sur $Q(Y)$.

\VerdierSemanticAnchor{ch2-2-3-4-b-c-e-proof}
L'assertion b) résulte de a), iii). L'assertion c) résulte de a), v). Il
résulte aussi de a), v) que les objets de $\mathcal D/\mathcal B$ qui sont
libérables à droite sont les objets de $\mathcal D/\mathcal B$ qui sont
isomorphes dans $\mathcal D/\mathcal B$ aux images par $Q$ des objets
$Q$-libres à droite. Les flèches de liberté sont des isomorphismes. Ceci
démontre e). Il est alors clair que $\operatorname{Lib}_R(Q)$ est une
sous-catégorie triangulée strictement pleine.

\VerdierSemanticAnchor{ch2-2-3-4-d-proof}
Soit $Z=Z_1\oplus Z_2$ un objet de $\operatorname{Lib}_R(Q)$. Il existe donc
un objet $Q$-libre à droite $Y$ et un isomorphisme
$Q(Y)\xrightarrow{\sim}Z_1\oplus Z_2$. Le foncteur
$X\longmapsto\operatorname{Hom}_{\mathcal D}(X,Y)$ admet donc un projecteur,
et par suite la décomposition de $Z$ en somme directe définit sur $Y$ un
projecteur. Les facteurs de $Z$ seront dans $\operatorname{Lib}_R(Q)$ si et
seulement si ce projecteur est décomposable, ce qui est le cas lorsque dans
$\mathcal D$ les sommes ou produits dénombrables sont représentables
(1.2.9).
\endgroup
